%
% Bachelor thesis: Autonomous Detection and Correction of Mixing Problems Using an AI Agent
%
% Based on the LaTeX template of the Institute for Optical Systems (IOS), HTWG Konstanz
% (Nico Brügel, adjusted by Lorenz Bung, 01.07.2020).
%
% Build (from this folder):  latexmk -xelatex thesis.tex
%

% Include header
% !TEX root = ../thesis.tex
\documentclass[11pt,a4paper,oneside]{report}		                     

% Mathe
\usepackage{marvosym}		% Package für Euro&Co.
\usepackage{amsfonts}		% Package für Mengensymbole (IN, IR, ...)
\usepackage{amsmath}		% Package für restlichen Mathekram
\usepackage{amssymb}        % math. Symbole und Umgebungen
\usepackage{graphicx}
\usepackage{mathtools}
\usepackage{unicode-math}   % Erlaubt die Mathe Schrift zu setzen
\newtheorem{definition}{Definition}		% Umgebung für Definitionen
\newcommand{\definitionautorefname}{Definition}		% Autoref für Definitionen

\usepackage{algorithm}		% Pseudocode listings
\usepackage{algpseudocode}
\algrenewcommand{\algorithmiccomment}[1]{\hskip3em$//$ #1}	% Kommentarformat
\newcommand{\algorithmautorefname}{Algorithmus}		% Autoref für Pseudocode listings

\usepackage{caption}
% Customizing figure numbering to make it bold
\captionsetup[figure]{labelfont=bf}
% Use german umlaute
%\usepackage{ngerman}
% \usepackage[ngerman]{babel}
%\usepackage[T1]{fontenc}
%\usepackage[utf8]{inputenc}
\usepackage[english]{babel}

% Zitierung und Literaturverzeichnis
%\usepackage{cite}
%\usepackage{apacite}		% Zitierstil
%\usepackage{natbib} 		% Erweiterte zitiermöglichkeiten 
%\usepackage{babelbib}		% Übersetzung des Literaturverzeichnisses und der Zitate
\usepackage[%
%,style=authoryear    % Alphabeticalsch
 ,style=ieee  % numerical-compressed
 ,citestyle=numeric-comp
 ,sorting=none        % no sorting
 ,sortcites=true      % some other example options ...
 ,block=none
 ,indexing=false
 ,citereset=none
 ,isbn=true
 ,url=true
 ,doi=true            % prints doi
 ,natbib=true         % if you need natbib functions
,backend=biber]{biblatex}
\addbibresource{bib/references.bib}

% Fonts
\usepackage{fontspec}
% \setmainfont[Path=./fonts/, BoldFont=swis721-Bold.ttf, ItalicFont=swis721-Italic.ttf, BoldItalicFont=swis721-Bold-Italic.ttf]{swis721.ttf}
% \setmonofont[Path=./fonts/, BoldFont=CourierNew-Bold.ttf, ItalicFont=CourierNew-Italic.ttf, BoldItalicFont=CourierNew-Bold-Italic.ttf]{CourierNew.ttf}
% \setmathfont[Path=./fonts/]{CambriaMath.ttf}


% formatting
\usepackage{a4}

% PDF als Cover Hintergrund
\usepackage[firstpage=true]{background}
\usepackage{geometry}
\usepackage{tikz}
\usepackage[export]{adjustbox} %outline around images


%%%%%%%%%%%%
%% HTWG CI Colors
%%%%%%%%%%%%
\definecolor{htwg-white}{cmyk}{0.0,0.0,0.0,0.0}
\definecolor{htwg-black}{cmyk}{0.0,0.0,0.0,1.0}
\definecolor{htwg-black-90}{cmyk}{0.0,0.0,0.0,0.9}
\definecolor{htwg-black-80}{cmyk}{0.0,0.0,0.0,0.8}
\definecolor{htwg-black-70}{cmyk}{0.0,0.0,0.0,0.7}
\definecolor{htwg-black-60}{cmyk}{0.0,0.0,0.0,0.6}
\definecolor{htwg-black-50}{cmyk}{0.0,0.0,0.0,0.5}
\definecolor{htwg-black-40}{cmyk}{0.0,0.0,0.0,0.4}
\definecolor{htwg-black-30}{cmyk}{0.0,0.0,0.0,0.3}
\definecolor{htwg-black-20}{cmyk}{0.0,0.0,0.0,0.2}
\definecolor{htwg-black-10}{cmyk}{0.0,0.0,0.0,0.1}
\definecolor{htwg-soft-blue}{cmyk}{0.1,0.0,0.0,0.1}
\definecolor{htwg-dark-blue}{rgb}{0.165, 0.22, 0.29}
\definecolor{htwg-teal}{cmyk}{1.0,0.0,0.5,0.0}

% Listing colours (used by \lstset below)
\definecolor{lighter-gray}{gray}{0.96}
\definecolor{ll-gray}{gray}{0.85}
\definecolor{l-gray}{gray}{0.7}

% Visible writing markers; remove all \todo{} before submission
\newcommand{\todo}[1]{\textcolor{red}{\textbf{[TODO: #1]}}}


% Verbesserte Darstellug der Schriftzeichen und Zwischenräumen
%\usepackage{lmodern}				% Benutze skalierbare Schriftfamilie
%\usepackage[activate]{microtype}

% Heading styles
\usepackage[sf,bf,raggedright]{titlesec}
\titlespacing*{\section}
{0pt}{5.5ex plus 1ex minus .2ex}{4.3ex plus .2ex}
\titlespacing*{\subsection}
{0pt}{5.5ex plus 1ex minus .2ex}{4.3ex plus .2ex}
\titleformat{\chapter}[hang]{\Huge\bfseries}{\thechapter.}{15pt}{\Huge\bfseries\color{htwg-teal}}
\titleformat{\section}[hang]{\Large\bfseries}{\thesection}{15pt}{\Large\bfseries\color{htwg-teal}}
\titleformat{\subsection}{\large\bfseries}{\thesubsection}{15pt}{\large\bfseries\color{htwg-teal}}
\titleformat{\subsubsection}[hang]{\normalsize\bfseries}{\thesubsubsection}{15pt}{\normalsize\bfseries\color{htwg-teal}}

% Header and Footer
\usepackage{fancyhdr}

\fancypagestyle{fancybody}{
	\fancyhf{}
	\fancyhead[L]{\nouppercase{\leftmark}}
	\fancyfoot[C]{\thepage}
	\renewcommand{\headrule}{\color{htwg-teal}\hrule}
}

\fancypagestyle{plain}{
	\fancyhf{}
	\fancyfoot[C]{\thepage}
	\renewcommand{\headrule}{}
}

\pagestyle{fancybody}
\renewcommand{\chaptermark}[1]{\markboth{\normalfont\thechapter.\space\color{htwg-teal}#1}{}}

% Absatz einrücken unterdrücken
\setlength{\parindent}{0pt}

% Listings
\usepackage{listings}
\lstset{
	backgroundcolor=\color{lighter-gray},
	xleftmargin=\parindent,
	basicstyle=\ttfamily\footnotesize,
	breaklines=true,
	captionpos=b,
	showspaces=false,
	showstringspaces=false,
	frame=shadowbox,
	rulesepcolor=\color{ll-gray},
	rulecolor=\color{l-gray},
	belowskip=11pt,
}

% Depth of chapters and subsections
\setcounter{tocdepth}{3}
\setcounter{secnumdepth}{3}

% Url line break
\def\UrlBreaks{\do\/\do-}
\def\UrlBigBreaks{\do\:\do\/}

% Links
\usepackage[pdfstartview=FitV, 	% start in 'fit size height'-view
colorlinks=true, 		% Links farbig markieren
linkcolor=htwg-teal, 	% Interne Links
citecolor=htwg-teal, 	% Links zur Literatur
filecolor=black, 		% Links auf lokale Dateien
urlcolor=htwg-teal,   	% Externe Links blau
breaklinks=true, 		% Links umbrechen
linktoc=page 	% Im Inhaltsverzeichnis sind Seitenzahlen UND Text Links
]{hyperref}

%Move to the right.
\setlength{\hoffset}{5pt}
\setlength{\parskip}{\baselineskip}

% Suppress warning of too small headheight
\headheight13.6pt

% Abkuerzungsverzeichnis
\usepackage[printonlyused]{acronym}
\usepackage[mode=buildnew]{standalone}

\usepackage[toc, page]{appendix}


% Personal data (name, student ID, e-mail, birth date and place) lives in
% local/personal.tex, which is gitignored because the repository is public.
% Copy local/personal.example.tex to local/personal.tex and fill it in.
\InputIfFileExists{local/personal}{}{% !TEX root = ../thesis.tex
% Template for personal data. Copy this file to local/personal.tex (gitignored)
% and replace the placeholders. Without local/personal.tex the build uses this file.
\newcommand{\autor}{\todo{author name}}
\newcommand{\matrikelnummer}{\todo{student ID}}
\newcommand{\autorEmail}{\todo{e-mail}}
\newcommand{\autorGeburtsort}{\todo{place of birth}}
\newcommand{\autorGeburtsdatum}{\todo{date of birth}}
}

% Thesis metadata
\newcommand{\thema}{Autonomous Detection and Correction of Mixing Problems Using an AI Agent}
\newcommand{\arbeitstyp}{Bachelor Thesis}
\newcommand{\abgabedatum}{18.01.2027}
\newcommand{\studiengang}{Applied Computer Science}
\newcommand{\prueferA}{\todo{first examiner}}
\newcommand{\prueferB}{\todo{second examiner}}


\begin{document}

% Title pages
% !TEX root = ../thesis.tex
\begin{titlepage}

\clearpage
\newgeometry{margin=0pt}
\thispagestyle{empty}

%% We need the this to perform coordinate calculations in TikZ.
\usetikzlibrary{calc}%
\usetikzlibrary{backgrounds}
\begin{tikzpicture}[remember picture,overlay, background rectangle/.style={fill=htwg-dark-blue}, show background rectangle]

% HTWG design elements (raster, lines etc)
\begingroup
    \foreach \x in {1,...,9} {
        \foreach \y in {1,...,13} {
        \pgfmathtruncatemacro{\xpos}{20.8 * \x}
        \pgfmathtruncatemacro{\ypos}{-20.8 * \y}
        \node (x\x y\y) at (\xpos mm, \ypos mm)[anchor=center,circle,draw,htwg-white,fill=htwg-white,scale=0.1]{};
        }
    }
    \draw [htwg-white] (-2mm, -115mm) -- (41mm, -166mm);
    \draw [htwg-white] (41mm, -270mm) -- (145mm, -228mm);

    \draw [line width=15mm,cap=round,htwg-teal] (62mm, -124mm) -- (187mm, -166mm);

    \node at (83mm, -124mm)[anchor=center,circle,draw,htwg-white,fill=htwg-white,scale=2.0]{};
    \node at (145mm, -228mm)[anchor=center,circle,draw,htwg-white,fill=htwg-white,scale=1.5]{};
\endgroup

% BSC letters
\begingroup
    \node at (15.8mm,-82.1mm)[anchor=center, fill=htwg-dark-blue]{
        \includegraphics[scale=0.165]{fonts/B_black}
    };
    \node at (36.8mm,-250.1mm)[anchor=center]{
        \includegraphics[scale=0.17]{fonts/S_black}
    };
    \node at (162.8mm,-124.1mm)[anchor=center]{
        \includegraphics[scale=0.15]{fonts/C_black}
    };
\endgroup

% HTWG-IN and IOS logos
\begingroup
    \node at (16.5mm, -15mm)[anchor=north west, fill=htwg-dark-blue]{
            \begin{minipage}[t]{30mm}
                \includegraphics[width=26mm]{images/htwg_in.png}
            \end{minipage}
        };
    \endgroup


% Thesis title
\node at (17.5mm,-180.5mm)[anchor=north west, fill=htwg-dark-blue]{
    \begin{minipage}[t]{145mm}
        \begin{flushleft}
            \color{htwg-white}\large{\arbeitstyp}\\[4mm]
            \color{htwg-white}\Huge{\thema}
        \end{flushleft}
    \end{minipage}
};

% Faculty etc
\node at (134.8mm,-17mm)[anchor=north west, fill=htwg-dark-blue]{
    \begin{minipage}[t]{55mm}
        \begin{flushleft}
            \color{htwg-white}\scriptsize\textbf{University of Applied Sciences Constance}
            \linebreak
            \color{htwg-white}\scriptsize{Faculty of Computer Science}
        \end{flushleft}
    \end{minipage}
};


% Author information
\node at (134.8mm,-38mm)[anchor=north west, fill=htwg-dark-blue]{
    \begin{minipage}[t]{55mm}
        \begin{flushleft}
            \color{htwg-white}\scriptsize\textbf{Submitted by}
            \linebreak
            \color{htwg-white}\scriptsize{\autor}
            \linebreak 
            \color{htwg-white}\scriptsize{Student ID \matrikelnummer}
            \linebreak
            \color{htwg-white}\scriptsize{\autorEmail}
        \end{flushleft}
    \end{minipage}
};

% Examiners and degree programme
\node at (134.8mm,-62mm)[anchor=north west, fill=htwg-dark-blue]{
    \begin{minipage}[t]{55mm}
        \begin{flushleft}
            \color{htwg-white}\scriptsize\textbf{Examiners}
            \linebreak
            \color{htwg-white}\scriptsize{\prueferA}
            \linebreak
            \color{htwg-white}\scriptsize{\prueferB}
            \linebreak
            \color{htwg-white}\scriptsize{\studiengang}
        \end{flushleft}
    \end{minipage}
};


\end{tikzpicture}
%% Restore the margins and turn the page white again.
\restoregeometry
\pagecolor{white}
\end{titlepage}


\pagenumbering{Roman}
% !TEX root = ../thesis.tex
% TODO before submission: replace this text with the faculty's official declaration
% for bachelor theses (wording, language, and whether a declaration on the use of
% AI tools is required). The text below is a placeholder adapted from the template.
\thispagestyle{plain}
\chapter*{Affidavit}
\addcontentsline{toc}{chapter}{Affidavit}
\label{ch:affidavit}

\todo{Replace with the official HTWG declaration for bachelor theses, including any required statement on the use of AI tools.}

I hereby declare that I,
\textit{\autor, born on \autorGeburtsdatum\ in \autorGeburtsort},

\begin{tabular}{lp{12cm}}
(1) & have prepared this bachelor thesis independently under the supervision of \prueferA\ and have used no aids other than those specified;\\[1em]
(2) & have properly indicated the use of verbatim quotations, tables, drawings, images and programs from literature or other sources (including the Internet), as well as the use of ideas of other authors, at the appropriate places in the thesis.\\[1em]
\end{tabular}

\vspace*{1cm}

\noindent
I am aware that a false declaration will have legal consequences.\\

\vspace*{1cm}

\noindent
Constance, \abgabedatum \hfill \begin{tabular}{c} \\ \\ \rule{5cm}{1pt} \\ (Signature)\end{tabular}

\newpage

% !TEX root = ../thesis.tex
% Written last, once results exist. Keep to about 200-300 words: problem,
% approach, evaluation design, main findings (including negative ones), limitations.
\thispagestyle{plain}
\chapter*{Abstract}
\addcontentsline{toc}{chapter}{Abstract}
\label{ch:abstract}

\todo{Write after the evaluation (roadmap M8). Do not state findings before they exist.}

\newpage


\setlength{\parskip}{0pt plus1pt}	% no paragraph spacing in the table of contents
\tableofcontents
\newpage
\setlength{\parskip}{\baselineskip}

\renewcommand{\headheight}{14pt}

% !TEX root = ../thesis.tex
% Only acronyms actually used via \ac{...} are printed (printonlyused).
% Long forms are lower case so they read correctly when expanded mid-sentence;
% use \Ac{...} at the start of a sentence.
\thispagestyle{plain}
\chapter*{Acronyms}
\addcontentsline{toc}{chapter}{Acronyms}
\label{ch:acronyms}

\begin{acronym}[MAPE-K]
	\acro{ADR}{architecture decision record}
	\acro{AI}{artificial intelligence}
	\acro{API}{application programming interface}
	\acro{CLI}{command-line interface}
	\acro{DAW}{digital audio workstation}
	\acro{DSP}{digital signal processing}
	\acro{LLM}{large language model}
	\acro{MAPE-K}{monitor, analyse, plan, execute over shared knowledge}
	\acro{STFT}{short-time Fourier transform}
	\acro{WAV}{Waveform Audio File Format}
\end{acronym}

\newpage


% Keep roman page numbering continuous for the back matter
\newcounter{anhang}
\setcounter{anhang}{\value{page}}

\pagenumbering{arabic}

% Chapters
% ---------------------------------------------------------------------------
% !TEX root = ../thesis.tex
% Sources: docs/research/project-scope.md (v1.0), docs/project/requirements.md (v1.1),
% docs/research/methodology.md, ADR-004, ADR-006, ADR-008, ADR-010, ADR-027.
% Statements about the literature carry \todo{cite} until a verified source is added.
\chapter{Introduction}
\label{ch:introduction}

\section{Motivation}
\label{sec:intro-motivation}

Mixing turns a set of separately recorded or produced tracks into a single coherent piece of music. Alongside its creative side, mixing involves technical problems that can be described in terms of the audio signal: tracks that compete for the same frequency range at the same time, energy that accumulates in narrow frequency regions, or stereo content that partly cancels when the mix is reproduced in mono \todo{cite a mixing textbook or survey}. Finding and correcting such problems takes time and experience, which makes them a natural target for automation \todo{cite work on automatic or intelligent mixing}.

\Acp{LLM} are increasingly used as decision-making components that select actions from structured information and call external tools \todo{cite work on LLM agents and tool use}. This raises the question of whether such a component can decide how to correct technical mixing problems when it is given deterministic measurements of the audio, the declared role of each track and the producer's intent. At the same time, the outputs of an \ac{LLM} are not guaranteed to be correct or reproducible \todo{cite}. Letting an \ac{LLM} change audio directly would make its errors hard to detect and its results hard to attribute.

This thesis therefore treats the \ac{AI} agent as one component of a larger system rather than as the system itself. The agent only proposes corrections. A deterministic control loop decides whether a proposal is admitted, renders the correction without modifying the source audio, measures the result again and either accepts the change or rolls it back. Whether this combination actually improves technical mixing problems, and under which conditions, is the subject of the evaluation.

\section{Problem Statement}
\label{sec:intro-problem}

Three difficulties make the task more than an engineering exercise.

First, the problems must be measurable. A deterministic measurement can show, for example, that two tracks contain energy in the same frequency band at the same time. Such overlap is at most an indication of possible masking; it is not proof that a listener perceives one track as masked by the other \todo{cite psychoacoustic masking literature}. The thesis keeps this distinction explicit throughout: measurement overlap, spectral peaks or phase correlation are never presented as proof of an audible problem.

Second, corrections must be safe. A proposed correction may be outside sensible parameter ranges, conflict with the producer's intent or make other parts of the mix worse. The system must reject such proposals before they are executed, detect regressions after a correction and restore the previously accepted state exactly.

Third, the evaluation must be fair and must not leak the answer. If the decision process or the runtime checks could see the clean reference audio or the parameters of an injected defect, they would have information that is not available in real use, and the results would be inflated. The agent-guided approach must also be compared with a credible deterministic baseline under otherwise equal conditions, so that observed differences can be attributed to the decision process.

\section{Research Questions and Hypotheses}
\label{sec:intro-rq}

The main research question of this thesis is:

\begin{quote}
\textit{To what extent can an AI agent autonomously detect, correct and verify selected technical mixing problems in multitrack music using deterministic audio analysis while respecting producer intent?}
\end{quote}

It is divided into four subquestions:

\begin{description}
	\item[RQ1] How can selected mixing problems be represented as measurable observations?
	\item[RQ2] Can an \ac{LLM} agent choose appropriate actions from measurements, context and producer intent?
	\item[RQ3] Does iterative correction outperform an otherwise comparable one-pass approach?
	\item[RQ4] Can validation and rollback prevent unsafe or unsuccessful interventions?
\end{description}

The following working hypotheses are tested, not assumed:

\begin{description}
	\item[H1] On controlled cases, agent-guided corrections improve predetermined technical objectives.
	\item[H2] Re-analysis and revision increase success or reduce unresolved defects compared with a single pass.
	\item[H3] Constraints and verification reduce undesirable or out-of-intent interventions.
\end{description}

Negative, partial and failed outcomes are reported as results in their own right.

\section{Scope}
\label{sec:intro-scope}

The thesis addresses three families of technical mixing problems:

\begin{description}
	\item[P1] spectral and temporal masking between tracks,
	\item[P2] frequency buildup and resonance,
	\item[P3] stereo and phase compatibility.
\end{description}

The interaction between kick drum and bass is the primary use case for P1, but it is not meant to be an assumption built into the system: the architecture is designed for any number of tracks with declared roles. \todo{State how far the final implementation meets this; the current analyser still takes exactly two stems labelled kick and bass, a known deviation under review (ADR-006).} All three families are research targets. Their measurement methods and implementation effort are expected to differ and are defined per family after a literature review. As a contingency, the minimum scientifically meaningful result is P1 implemented end-to-end and evaluated under all comparison conditions; the depth reached for P2 and P3 is reported as achieved.

The system works offline on synchronised \ac{WAV} stems exported from multitrack projects, with track roles declared manually. Audio analysis and corrections are deterministic \ac{DSP} written in Python. The decision agent uses the Claude \ac{API} and returns structured proposals. Source audio is never modified, and every correction is reversible.

Outside the scope are the training of new models, comprehensive automatic mixing or mastering, real-time or plug-in operation, a graphical user interface, direct manipulation of \ac{DAW} project files, systems with several cooperating agents, and any claim of guaranteed perceptual superiority. The study material covers house and techno, trap, phonk and wave, and drum and bass; comparisons between genres are exploratory unless the data supports them.

\section{Approach}
\label{sec:intro-approach}

The prototype is a feedback control loop inspired by the \ac{MAPE-K} reference model \todo{cite Kephart and Chess or a self-adaptive systems source}. Each iteration measures the stems, identifies plausible problem candidates together with the strength of the evidence for them, asks the decision agent for a proposal, admits or rejects that proposal, renders the admitted correction, and verifies whether the targeted problem improved without making other measurements worse. Depending on the outcome, the loop accepts the new state or rolls back to the last accepted one. Iterations, actions, \ac{AI} requests and runtime are bounded, and every stop has a recorded reason.

The decision agent cannot access files, execute processing or bypass admission. All analysis, admission, rendering and verification are deterministic and can be tested without network access. The software is organised as a hexagonal modular monolith: the control loop and its rules form the core, while audio analysis, rendering and the connection to Claude are adapters behind interfaces defined by the core.

The evaluation compares four conditions on the same material: the unprocessed input (B0), a deterministic rule-based correction (B1), a single pass guided by the agent (B2) and iterative correction guided by the agent (B3). The test material consists of real stems from self-created projects into which defects are injected in a controlled and documented way. The clean reference audio and the injection parameters are stored separately and are used only after a run, for research evaluation. Thresholds and parameters are chosen on a calibration split and frozen before the evaluation split is used. Because the agent's outputs vary between calls, the agent-guided conditions are repeated, and the spread of results is reported.

\section{Contributions}
\label{sec:intro-contributions}

The thesis aims to contribute:

\begin{itemize}
	\item a measurable representation of selected technical mixing problems, with an explicit statement of evidence strength and of the gap between measurement and perception;
	\item the design and implementation of a control loop in which an \ac{AI} agent proposes corrections and deterministic components admit, render, verify and, if needed, roll them back;
	\item an evaluation design that separates runtime information from hidden ground truth and compares the agent with a deterministic baseline under declared, otherwise equal conditions;
	\item an empirical assessment of H1--H3 on controlled cases, including negative and failed outcomes.
\end{itemize}

\todo{Revise this list once the results are known; state only contributions that were delivered.}

\section{Structure of the Thesis}
\label{sec:intro-structure}

\autoref{ch:background} introduces the technical mixing problems, prior work on automatic mixing and the use of \acsp{LLM} as decision makers. \autoref{ch:requirements} derives the requirements and the research design from the research questions. \autoref{ch:architecture} presents the architecture of the system and \autoref{ch:implementation} its implementation. \autoref{ch:evaluation} describes the dataset, the comparison conditions and the measures, and \autoref{ch:results} reports the results. \autoref{ch:discussion} answers the research questions and discusses limitations, and \autoref{ch:conclusion} concludes with an outlook.

% !TEX root = ../thesis.tex
\chapter{Background and Related Work}
\label{ch:background}
% Sources: primary literature only (no citation without a verified source).
% The roadmap M1 literature review feeds this chapter; see the open questions in
% docs/research/project-scope.md.

\section{Multitrack Mixing and Technical Mixing Problems}
\label{sec:bg-mixing}
\todo{Role of mixing; what technical (as opposed to aesthetic) problems are.}

\subsection{Spectral and Temporal Masking (P1)}
\label{sec:bg-masking}
\todo{Psychoacoustic masking; multitrack masking measures from the literature; difference between time--frequency overlap and perceived masking (SR-08).}

\subsection{Frequency Buildup and Resonance (P2)}
\label{sec:bg-buildup}
\todo{Definitions and published detection methods.}

\subsection{Stereo and Phase Compatibility (P3)}
\label{sec:bg-stereo}
\todo{Phase correlation and mono compatibility; published methods.}

\section{Automatic and Intelligent Mixing}
\label{sec:bg-automix}
\todo{Rule-based, knowledge-engineered and learned approaches to automatic mixing, and how they are evaluated.}

\section{Large Language Models as Decision Makers}
\label{sec:bg-llm}
\todo{LLM agents and tool use; structured outputs; variability of LLM outputs (motivates SR-05).}

\section{Feedback Control Loops in Self-Adaptive Systems}
\label{sec:bg-mapek}
\todo{The MAPE-K reference model, used to structure the control loop (ADR-027).}

\section{Research Gap}
\label{sec:bg-gap}
\todo{What prior work leaves open that this thesis addresses; claim only what the cited sources support.}

% !TEX root = ../thesis.tex
\chapter{Requirements and Research Design}
\label{ch:requirements}
% Sources: docs/project/requirements.md (v1.1), docs/research/project-scope.md,
% docs/research/methodology.md, ADR-006, ADR-008, ADR-010.

\section{Operationalising the Research Questions}
\label{sec:req-rq}
\todo{How RQ1--RQ4 and H1--H3 from \autoref{sec:intro-rq} are turned into measurable criteria.}

\section{System Context and Scope}
\label{sec:req-context}
\todo{Inputs, outputs, external service, operator and exclusions (requirements.md, section 2).}

\section{Functional Requirements}
\label{sec:req-functional}
\todo{Summary of FR-01 to FR-18 (FR-14 withdrawn) and their traceability to the research questions.}

\section{Quality, Scientific and Operational Requirements}
\label{sec:req-quality}
\todo{QR, SR and OP requirements, in particular determinism (QR-02), the hidden-reference boundary (SR-10) and non-circularity (SR-03).}

\section{Comparison Conditions}
\label{sec:req-conditions}
\todo{B0 unprocessed, B1 rule-based, B2 agent single-pass, B3 agent iterative; independent variables and controlled factors per comparison (SR-01). Depends on open decision D7.}

\section{Open Design Decisions}
\label{sec:req-open}
\todo{Decisions open at the time of writing (D2--D11) and how each was resolved; update as they are decided.}

% !TEX root = ../thesis.tex
% Sources: docs/architecture/overview.md, docs/project/engineering-conventions.md,
% ADR-001, ADR-004, ADR-005, ADR-010, ADR-023, ADR-026 and ADR-027 (all in docs/decisions/).
% State of the code described here: 2026-11-04 (ADR-027 migration completed, no control loop yet).
% Update the status table and the \todo markers when the loop, diagnosis and decision agent exist.
\chapter{System Architecture}
\label{ch:architecture}

This chapter describes the architecture of the prototype: which requirements shape it, how the code is divided into parts, how those parts interact in the control loop and where the boundaries lie that the research design depends on. Architectural choices are recorded as \acp{ADR} in the project repository \cite{AutonomousMixingSystem}; the chapter refers to them by number (for example ADR-027).

\section{Architectural Drivers}
\label{sec:arch-drivers}

Most of the requirements in \autoref{ch:requirements} concern the behaviour of single components. A few of them, however, constrain how the system as a whole must be organised, and these drive the architecture:

\begin{description}
	\item[Attributable decisions.] Every change to the audio must be traceable to a proposal, a decision to admit it and a verification result (FR-07, FR-12). The \ac{AI} component therefore only proposes; it cannot access files, run processing or bypass admission (OP-01, ADR-004).
	\item[Offline determinism.] Analysis, admission, rendering and verification must be deterministic and testable without network access or credentials (QR-02), and must not depend on a particular \ac{AI} provider (QR-07).
	\item[Non-destructive, reversible changes.] Source audio is never modified (QR-01), a run never overwrites the outputs of an earlier run (OP-02), and a rejected change must restore the last accepted state exactly (FR-10).
	\item[Hidden-reference boundary.] The clean reference audio and the ground truth of injected defects must be unreachable for every controller, its decision inputs and runtime verification (SR-10, ADR-010).
	\item[Exchangeable controllers.] The comparison conditions B0--B3 must run on the same inputs, measurements and action capabilities and be selected by configuration only (SR-01, FR-15).
\end{description}

The system is a research prototype operated by one researcher through a \ac{CLI} and runs offline on \ac{WAV} files (ADR-001). It is therefore built as a single deployable Python package, a \emph{modular monolith}, rather than as distributed services.

\section{Hexagonal Modular Monolith}
\label{sec:arch-hexagonal}

The code follows the hexagonal or ports-and-adapters style \cite{Cockburn2005}. Its central idea is that the application core defines the interfaces (\emph{ports}) through which it talks to the outside world, and that everything with side effects, such as file access, signal processing or calls to an external service, is implemented as an \emph{adapter} behind one of these ports. The core can then be run and tested with stand-ins for every adapter.

The package \texttt{autonomous\_mixing} is divided into four rings (ADR-027):

\begin{description}
	\item[\texttt{core}] contains the control loop, the pure rules it applies at each stage and the immutable records these rules exchange. It imports nothing outside itself: no NumPy, no audio library, no provider SDK and no file system access.
	\item[\texttt{adapters}] implement the ports: deterministic audio analysis (\texttt{dsp\_analysis}), rendering of corrections including \ac{WAV} input and output (\texttt{rendering}), the Claude-based decision component (\texttt{decision\_agent}) and alternative decision makers for the baseline conditions (\texttt{baselines}).
	\item[\texttt{drivers}] build the adapters and start runs: the \ac{CLI} and, later, the experiment runner (\texttt{experiments}).
	\item[\texttt{evaluation}] holds the hidden reference, the ground truth of injected defects and the research metrics, and is used only after a run.
\end{description}

Inside the core, \texttt{records} defines all data types, \texttt{ports} the three port protocols and the checks on their responses, \texttt{admission} and \texttt{verification} the rules of the corresponding stages, \texttt{diagnosis} the identification of problem candidates (from its roadmap phase on), and \texttt{loop} the composition of all stages. \autoref{tab:arch-packages} lists the packages and their state at the time of writing.

\begin{table}[htb]
	\centering
	\small
	\begin{tabular}{@{}>{\raggedright\arraybackslash}p{0.31\textwidth}>{\raggedright\arraybackslash}p{0.42\textwidth}>{\raggedright\arraybackslash}p{0.21\textwidth}@{}}
		\toprule
		Package & Responsibility & State \\
		\midrule
		\texttt{core.records} & immutable records and their invariant checks & implemented \\
		\texttt{core.ports} & \texttt{AudioAnalyzer}, \texttt{DecisionMaker}, \texttt{Renderer}; response checks & implemented \\
		\texttt{core.admission} & admission of proposals; metadata preflight & implemented \\
		\texttt{core.verification} & accept-or-roll-back rules & record checks only \\
		\texttt{core.diagnosis} & problem candidates and evidence strength & planned \\
		\texttt{core.loop} & stage order, budgets, transitions, rollback & state-transition checks only \\
		\texttt{adapters.dsp\_analysis} & deterministic measurements & descriptive analysis (M0) \\
		\texttt{adapters.rendering} & gain and mix rendering, \ac{WAV} I/O, integrity checks & implemented \\
		\texttt{adapters.decision\_agent} & Claude-based \texttt{DecisionMaker} & planned \\
		\texttt{adapters.baselines} & other \texttt{DecisionMaker}s (B1) & planned \\
		\texttt{drivers.cli} & command-line entry point & implemented \\
		\texttt{drivers.experiments} & runs all conditions over a data split & planned \\
		\texttt{evaluation} & hidden reference and research metrics & planned \\
		\bottomrule
	\end{tabular}
	\caption{Packages of the prototype and their state on 4 November 2026.}
	\label{tab:arch-packages}
\end{table}

\todo{Update \autoref{tab:arch-packages} before submission.}

Four import rules make the boundaries checkable (ADR-027): (a) \texttt{core} imports nothing outside \texttt{core}; (b) inside the core, \texttt{records} imports nothing outside itself, \texttt{ports} imports only \texttt{records}, the stage packages import only \texttt{records} and \texttt{ports}, and only \texttt{loop} composes stages and ports; (c) adapters import only \texttt{core.records}, \texttt{core.ports}, third-party libraries and their own modules, never another adapter or the loop; (d) only \texttt{drivers.experiments} imports \texttt{evaluation}. Architecture tests in the repository parse the source of every module and fail when a rule is broken. These are static checks: imports constructed at run time are not analysed.

The previous layout divided the same code into horizontal domain, application and infrastructure layers. In that layout one concept, such as a decision, was spread over three packages, and no component connected the stages, so the system could not be read from end to end. The hexagonal layout replaced it on 4 November 2026 (ADR-027, superseding ADR-017 and ADR-022). The rules above express exactly the boundaries the thesis relies on and little more.

\section{The Control Loop}
\label{sec:arch-loop}

The behaviour of the core is a feedback loop modelled on the \ac{MAPE-K} reference model of autonomic computing \cite{Kephart2003}. In \ac{MAPE-K}, an autonomic manager \emph{monitors} a managed resource, \emph{analyses} the collected data, \emph{plans} a change and \emph{executes} it, all over shared \emph{knowledge}. Here the managed resource is the state of a mixing project, and the knowledge is the run history. \autoref{fig:arch-loop} shows the stages.

\begin{figure}[htb]
	\centering
	\begin{tikzpicture}[
		stage/.style={draw=htwg-dark-blue, rounded corners=2pt, minimum width=24mm, minimum height=9mm, align=center, font=\small},
		ai/.style={stage, fill=htwg-soft-blue, dashed},
		det/.style={stage, fill=white},
		flow/.style={-{Stealth[length=2mm]}, thick, htwg-dark-blue},
		note/.style={font=\scriptsize, align=center, text=htwg-black-70},
		node distance=6mm and 9mm
	]
		\node[det] (measure) {Measure};
		\node[det, right=of measure] (diagnose) {Diagnose};
		\node[ai, right=of diagnose] (decide) {Decide};
		\node[det, right=of decide] (admit) {Admit};
		\node[det, below=14mm of admit] (render) {Render};
		\node[det, left=of render] (verify) {Verify};
		\node[det, left=of verify] (accept) {Accept or\\ roll back};
		\node[det, left=of accept] (history) {Run history};

		\draw[flow] (measure) -- (diagnose);
		\draw[flow] (diagnose) -- (decide);
		\draw[flow] (decide) -- (admit);
		\draw[flow] (admit) -- node[note, right]{admitted\\ action} (render);
		\draw[flow] (render) -- (verify);
		\draw[flow] (verify) -- (accept);
		\draw[flow] (accept) -- (history);
		\draw[flow] (history) -- node[note, left]{next\\ iteration} (measure);
		\draw[flow, dashed] (admit.north) -- ++(0,5mm) -| node[note, above, pos=0.25]{rejected: recorded, no rendering} (history.north east);

		\node[note, below=6mm of verify] {verify re-measures the candidate and compares it with the last accepted state};
	\end{tikzpicture}
	\caption{Stages of the control loop. Only the dashed stage is performed by the \ac{AI} component; all other stages are deterministic. Iterations, actions, \ac{AI} requests and runtime are bounded, and every stop has a recorded reason.}
	\label{fig:arch-loop}
\end{figure}

The stages map to \ac{MAPE-K} as follows: \emph{measure} corresponds to Monitor and \emph{diagnose} to Analyze. Plan is split into two steps: the decision maker \emph{decides} by proposing one action or abstaining, and a deterministic \emph{admission} step decides whether that proposal may be executed. Execute corresponds to \emph{rendering} the admitted action into new candidate audio. Unlike the reference model, the loop then \emph{verifies} the candidate explicitly: it measures the candidate again and checks whether the targeted problem improved without making other measurements worse, at the level of individual stems and of the combined mix (FR-09). The loop accepts the candidate or rolls back to the last accepted state, records the attempt and continues with the next iteration until the decision maker abstains, a budget is exhausted or a failure ends the run.

The stage order is fixed in the code of \texttt{core.loop}. There is no event bus, plug-in mechanism or configurable pipeline: the loop calls adapters, and adapters never call the loop. This keeps a run readable from beginning to end and makes each stage testable with plain values.

The records for this loop already exist. A run record (\texttt{RunRecord}) holds the configuration, budgets and their usage, and the outcome of the run (succeeded, abstained, limit reached, failed or cancelled). Each correction attempt (\texttt{CorrectionAttempt}) holds the complete state before the change and, if rendering succeeded, the candidate state, and ends as accepted, rolled back, failed or unresolved. Rollback restores the complete saved state, not an inverse operation, so restoring a state never depends on undoing a correction numerically.

\todo{The function that runs the loop (\texttt{core.loop.run}) does not exist yet. Describe its implemented behaviour, including budget enforcement and stop reasons, once it does.}

\section{Ports and Records}
\label{sec:arch-ports}

The core defines three ports as Python protocols:

\begin{description}
	\item[\texttt{AudioAnalyzer}] receives an analysis request for one project state and returns observations, that is, deterministic measurements with their provenance.
	\item[\texttt{DecisionMaker}] receives a \texttt{Situation}, which bundles the project state, the musical context with declared track roles and producer intent, and the observations, and returns a decision: either a proposed action or an abstention with a reason.
	\item[\texttt{Renderer}] receives a render plan for one admitted action and returns a rendered candidate or a failure.
\end{description}

Each port returns expected failures as values rather than exceptions, so that, for example, a failed request to the \ac{AI} service becomes a recorded outcome instead of a crash (OP-04). For each port the core also defines a check of the response against the request. A decision, for instance, must refer to the project revision, context and observations it was given and must name the controller that was asked; anything else is rejected before admission.

All data exchanged between stages and ports are \emph{records} in \texttt{core.records}. Every record is a frozen data class that validates itself on construction, so that an invalid record cannot exist. Its fields hold only immutable values, and new versions are derived by copying rather than by changing an existing record. Audio samples never enter records: records refer to audio files by identifier, location, format and checksum, while sample arrays stay inside the adapters. The outcome of a stage is a union type, for example an admitted or a rejected action, and the caller must handle every case explicitly.

Immutability serves testability and traceability: a request cannot change after it was sent, so the history can record exactly what a stage received. It is a precondition, not a guarantee. Rollback still depends on the retained audio files being present and passing integrity checks, and the history shows what the decision maker saw only if the actual request and the raw response are stored (FR-12, QR-03).

\section{Separation of Proposal and Execution}
\label{sec:arch-proposal}

The central architectural decision of the thesis is that the \ac{AI} component proposes and deterministic code disposes (ADR-004). A proposal is not permission to act. The admission step checks it against a versioned admission policy for the project, in a fixed order: the proposal must refer to the current project state, the policy must belong to the project, the proposal must contain an action at all, the target track must not be protected, an explicit permission for the action on that track must exist, and the requested parameter must lie within that permission's bounds. Admission never clamps a value into range; an out-of-bounds proposal is rejected and the reason is recorded (FR-07). An abstention is a valid outcome, not an error.

At the time of writing, the only action type is setting the absolute gain of one track (\texttt{SetTrackGain}). Gain is always set relative to the fixed input stem of the track, never as a change applied to the latest output, so repeated corrections do not accumulate. A \emph{preflight} step then checks, on metadata alone, that the admitted action can be rendered: the output workspace, the identity of the new candidate, the inputs of the mix and the target gain must all be consistent. The renderer for this action (ADR-026) reads and checks the fixed input of the target track, applies the gain, sums all stems into the combined mix in a declared order without normalisation or limiting, and writes the new target stem and the new mix as new files into a freshly allocated output directory. It never overwrites an existing file.

\todo{Further action types per problem family depend on open decision D5; extend this section when they exist.}

Because every decision maker sits behind the same port, the comparison conditions differ only in which implementation is plugged in. The rule-based baseline (B1) and the agent-guided conditions (B2, B3) receive the same situation and their proposals pass the same admission, rendering and verification. The exact definition of each condition is still open (D7).

\section{Hidden-Reference Boundary}
\label{sec:arch-hidden}

The evaluation must not leak its answer into the runtime (SR-10). The architecture contributes to this in two ways. First, hidden evaluation data lives only in the \texttt{evaluation} package, and only the experiment runner may import it (rule (d)); the core, the adapters and the \ac{CLI} cannot. Second, adapters read files only from roots that the caller supplies explicitly (ADR-023), so an adapter has no implicit path to the hidden data.

Neither measure is sufficient on its own. An import rule is not a file system sandbox: an adapter could still open a file by path. SR-10 therefore additionally requires that runtime inputs, decision inputs and runtime configuration contain no fields or references that point to hidden data, and that this is tested.

\todo{The acceptance tests for SR-10 do not exist yet. Report them here once they do, or list their absence as a limitation in \autoref{ch:discussion}.}

\section{Summary}
\label{sec:arch-summary}

The prototype is a hexagonal modular monolith whose core is a \ac{MAPE-K}-style control loop over immutable records. Audio analysis, rendering and the \ac{AI} decision component are adapters behind three ports. The \ac{AI} component only proposes; deterministic admission, rendering and verification decide what happens to the audio, and a rejected change is undone by returning to the complete saved state rather than by reversing the correction. Import rules, checked by tests, keep the core free of side effects and keep hidden evaluation data away from the runtime. \autoref{ch:implementation} describes how these parts are implemented.

% !TEX root = ../thesis.tex
\chapter{Implementation}
\label{ch:implementation}
% Describe only what is implemented and tested at the time of writing; mark planned parts as planned.

\section{Technology and Tooling}
\label{sec:impl-tooling}
\todo{Python, uv, test and lint setup (ADR-001, ADR-002).}

\section{Input Validation and WAV Handling}
\label{sec:impl-input}
\todo{Project loading, compatibility checks and error behaviour.}

\section{Measurement}
\label{sec:impl-measurement}
\todo{Deterministic DSP analysis; parameters and their sources, or their status as assumptions.}

\section{Diagnosis}
\label{sec:impl-diagnosis}
\todo{Problem candidates and evidence strength (FR-05), once implemented.}

\section{Decision Agent}
\label{sec:impl-agent}
\todo{Claude adapter: requests, structured proposals, failure handling, recorded raw requests and responses.}

\section{Admission, Rendering and Verification}
\label{sec:impl-admission}
\todo{Admission rules, gain/mix rendering, workspaces and integrity checks, verification and rollback.}

\section{Testing Strategy}
\label{sec:impl-testing}
\todo{Unit, contract and architecture tests; synthetic signals only in unit tests (ADR-006).}

% !TEX root = ../thesis.tex
\chapter{Evaluation Methodology}
\label{ch:evaluation}
% Sources: docs/research/methodology.md, docs/research/dataset-plan.md, ADR-008, ADR-010, SR-01 to SR-10.
% Freeze the criteria in this chapter before the main evaluation (SR-02, SR-03).

\section{Dataset and Controlled Defect Injection}
\label{sec:eval-dataset}
\todo{Self-created projects, injection protocol, calibration and evaluation splits.}

\section{Hidden Evaluation Reference}
\label{sec:eval-hidden}
\todo{What is hidden from the runtime, how it is stored and when it is joined with run records.}

\section{Measures}
\label{sec:eval-measures}
\todo{Reported measures, which of them the decision process or verification can see, and loudness control (D9).}

\subsection{Directional masking case study}
\label{sec:eval-masking-case}
The first case study will use a kick and bass from a mix exported by the producer from FL Studio. Both estimates will use the same aligned stems and the same evaluation interval. The analysis will report kick as masker and bass as target, then bass as masker and kick as target. Each result will identify the direction, time interval, frequency region and measurement configuration. When the evidence for a direction is unavailable, that direction will be reported as unassessed. The proposed estimator follows Ronan's cross-adaptive masking calculation for a pair of sources~\cite[Sec.~7.2.4.2, Eqs.~7.11--7.13]{Ronan2018Thesis}. For each direction, the other stem provides the masking threshold. On the common frame and scale-factor band grid, the signed margin is the threshold-to-target-energy ratio in decibels. This calculation requires positive threshold and target energy. A zero-energy target is unassessable. Ronan's positive-margin aggregation is recorded separately from the signed margin.

The planned reproduction will establish the margin calculation with known intermediate values first. Reproduction of the threshold front end requires a documented MPEG psychoacoustic model and an independent intermediate reference. The analysis will document the front-end settings and any departures from the original method before the case study.

Analytical fixtures will first check the estimator against known inputs, including silence, non-overlapping activity and a controlled change in masker level with the target held fixed. These checks will test the numerical implementation. For the music case, the same measurement configuration will be used before and after a single bounded gain change. Any decision criterion will be set on calibration material and kept fixed for evaluation.

A separate, level-matched listening comparison will examine the audibility of each part and any perceived change in mix balance. Where feasible, the listening order will be varied and condition labels concealed. The objective estimates and the listening observations will be reported separately because objective masking measures have shown limited agreement with perceived masking in multitrack mixes~\cite{Hafezi2015}. Listening observations from a single mix will be treated as case-study evidence.

\section{Experimental Procedure}
\label{sec:eval-procedure}
\todo{Conditions B0--B3, repeated trials (D8), budgets and the frozen configuration.}

\section{Threats to Validity}
\label{sec:eval-threats}
\todo{Injected versus real defects, overlap versus perceived masking, LLM variability, confounders between conditions.}

% !TEX root = ../thesis.tex
\chapter{Results}
\label{ch:results}
% Report only measured results from archived runs (roadmap M7), including negative and failed outcomes (O2, SR-06).

\todo{Written after the main evaluation.}

% !TEX root = ../thesis.tex
\chapter{Discussion}
\label{ch:discussion}

\section{Answers to the Research Questions}
\label{sec:disc-rq}
\todo{RQ1--RQ4 and H1--H3, each against the evidence in \autoref{ch:results}.}

\section{Limitations}
\label{sec:disc-limitations}
\todo{Including the absence or scale of a listening assessment (D10).}

% !TEX root = ../thesis.tex
\chapter{Conclusion and Outlook}
\label{ch:conclusion}

\section{Conclusion}
\label{sec:concl-conclusion}
\todo{Summary of what was built and what was found.}

\section{Outlook}
\label{sec:concl-outlook}
\todo{Future work, for example problem families not reached in depth, listening studies, DAW integration (deferred, ADR-009).}


% ---------------------------------------------------------------------------
\newpage
\pagenumbering{Roman}
\setcounter{page}{\value{anhang}}
\printbibliography[heading=bibintoc]
\listoffigures
\listoftables
% !TEX root = ../thesis.tex
\begin{appendices}

\chapter{Requirements Traceability}
\label{app:traceability}
\todo{Full traceability table from docs/project/requirements.md.}

\chapter{Reproducibility}
\label{app:reproducibility}
\todo{Repository version, frozen configuration identifier and commands to reproduce the evaluation.}

\end{appendices}


\end{document}
